% Gerado por scripts/etapa3_xgboost.py — requer \usepackage{booktabs}
\begin{table}[htbp]
\centering
\small
\caption{XGBoost — métricas de acurácia e de custo (sensibilidade: estoque arredondado para cima)}
\label{tab:xgb_metricas_sensibilidade}
\begin{tabular}{lrrrrrrrr}
\toprule
Quantis & MAE & RMSE & Pinball 80\% & Pinball 90\% & Pinball 95\% & Custo 1:4 & Custo 1:9 & Custo 1:19 \\
\midrule
Empírico & 1,246 & 2,281 & 0,532 & 0,381 & 0,251 & 2,727 & 3,838 & 4,985 \\
Normal & 1,246 & 2,281 & 0,540 & 0,379 & 0,253 & 2,847 & 3,836 & 4,918 \\
\bottomrule
\end{tabular}
\par\smallskip
{\footnotesize Nota: teste de 90 dias (26/01 a 24/04/2016), 5.748 séries. MAE e RMSE avaliam a previsão pontual e são iguais nos dois métodos de quantis. Pinball avalia o quantil previsto; o custo avalia a decisão de estoque. custo por série-dia com custo de sobra = 1; razão 1:4, 1:9 e 1:19 (sobra:falta) corresponde aos níveis de serviço de 80\%, 90\% e 95\%. Estoque = quantil arredondado para cima.}
\end{table}
